\section{Calibration sur le granite de Red Bohus}
\label{sec-calib}

Aucun jeu FDEM ne restitue à la fois la résistance en compression simple et l'enveloppe confinée du Red Bohus : les jeux calés sans confinement sous-estiment l'effet du confinement d'un facteur 2 à 3, et le seul jeu calé à 20~MPa surestime la résistance sous fort confinement de 25 à $36~\%$. L'émulateur de la campagne d'août a été infirmé par les calculs réels, et aucune calibration n'est valide aujourd'hui. Le jeu de juillet n'est plus reproductible avec le binaire courant (\cref{sec-calib-rejeu}).

Le Red Bohus est le matériau de la thèse, et aucune simulation d'impact ne portera sur lui tant que ses paramètres de joint ne seront pas calés. Trois campagnes FDEM ont été menées de juillet à septembre 2026, et des calibrations par lois continues sous Abaqus de mai à juillet.

\subsection{Cibles expérimentales}

Les essais de Dumoulin et al.~\cite{dumoulin2024data} portent sur un granite du sud-ouest de la Suède, à $60~\%$ de feldspath, $35~\%$ de quartz et $5~\%$ de biotite en masse (\cref{tab-cibles}).

\begin{table}[htbp]
\centering
\caption{Essais de Dumoulin et al. (2024) sur le granite de Red Bohus.}
\label{tab-cibles}
\small
\begin{tabular}{@{}lrll@{}}
\toprule
essai & nombre & grandeur & valeur \\
\midrule
compression simple & 4 & résistance au pic & $126{,}6 \pm 21{,}4$~MPa \\
brésilien & 4 & $2P/(\pi D t)$ & $10{,}27 \pm 0{,}98$~MPa \\
triaxial, $\sigma_3 = 20$~MPa & 3 & $\sigma_1$ au pic & $424{,}8$~MPa \\
triaxial, $\sigma_3 = 50$~MPa & 3 & $\sigma_1$ au pic & $649{,}0$~MPa \\
triaxial, $\sigma_3 = 75$~MPa & 3 & $\sigma_1$ au pic & $779{,}0$~MPa \\
triaxial, $\sigma_3 = 100$~MPa & 3 & $\sigma_1$ au pic & $899{,}3$~MPa \\
élasticité & 12 & $E$, $\nu$ & $77{,}7$~GPa, $0{,}29$ \\
\bottomrule
\end{tabular}
\end{table}

L'enveloppe est concave : sa pente $d\sigma_1/d\sigma_3$ passe de $14{,}9$ entre 0 et 20~MPa à $4{,}8$ entre 75 et 100~MPa. La dispersion de la compression simple atteint $17~\%$, contre $0{,}4$ à $1~\%$ sur les triaxiaux, et deux courbes sur quatre sont tronquées avant le pic. Trois jeux d'élasticité ont circulé ($52$~GPa et $0{,}25$ ; $57{,}3$~GPa et $0{,}17$ ; $77{,}7$~GPa et $0{,}29$), ainsi que deux cibles de traction, le brésilien mesuré de $10{,}3$~MPa et une traction directe équivalente de $18{,}3$~MPa tirée de la littérature. Un ajustement de Hoek et Brown,
\begin{equation}
\sigma_1 = \sigma_3 + \sigma_{ci}\sqrt{m_i\,\frac{\sigma_3}{\sigma_{ci}} + 1},
\end{equation}
donne $\sigma_{ci} = 178{,}5$~MPa et $m_i = 34{,}6$ ; il suit les essais confinés mais surestime la résistance simple de $41~\%$.

\subsection{Méthode bayésienne}

Les paramètres de joint $\boldsymbol\theta$ sont tirés d'une loi a priori uniforme sur une boîte $\Theta$ et confrontés aux cibles $\bar y_k$ d'écart-type $s_k$. Un calcul FDEM étant coûteux, chaque observable est remplacée par un émulateur gaussien $\eta_k$, ajusté sur un plan d'expériences, et l'écart structurel entre le modèle et la roche par un terme d'écart de modèle $\delta_k$, selon une forme simplifiée du cadre de Kennedy et O'Hagan où cet écart est réduit à une variance relative :
\begin{gather}
p(\boldsymbol\theta \mid \bar{\mathbf y}) \propto p(\bar{\mathbf y} \mid \boldsymbol\theta)\, \mathbb 1_\Theta(\boldsymbol\theta), \qquad
\bar y_k = \eta_k(\boldsymbol\theta) + \delta_k + \varepsilon_k, \\
\log p(\bar{\mathbf y} \mid \boldsymbol\theta) = -\frac12 \sum_k \left[\frac{\left(\mu_k(\boldsymbol\theta) - \bar y_k\right)^2}{v_k(\boldsymbol\theta)} + \ln v_k(\boldsymbol\theta)\right] + \text{cste}, \qquad
v_k(\boldsymbol\theta) = s_k^2 + \sigma_k^2(\boldsymbol\theta) + (d_k\,\bar y_k)^2,
\end{gather}
où $\mu_k$ et $\sigma_k^2$ sont la moyenne et la variance prédictives de l'émulateur et $d_k$ l'écart de modèle relatif (10 à $20~\%$). La variance $v_k$ dépend de $\boldsymbol\theta$ par l'émulateur, d'où le terme $\ln v_k$. Le code de juillet l'inclut ; celui d'août l'omet (\cle{calibration_redbohus/tools/analyze.py}), si bien que sa loi a posteriori favorise les régions où l'émulateur est le plus incertain, ce qui peut contribuer aux paramètres collés à leur borne. L'émulateur suit
\begin{equation}
\mu(\mathbf x_*) = m + \mathbf k_*^{\mathsf T}(\mathbf K + \sigma_n^2\mathbf I)^{-1}(\mathbf y - m), \qquad
\sigma^2(\mathbf x_*) = k_{**} - \mathbf k_*^{\mathsf T}(\mathbf K + \sigma_n^2\mathbf I)^{-1}\mathbf k_*,
\end{equation}
avec un noyau exponentiel carré en juillet et un noyau de Matérn $5/2$ à longueurs par dimension en août. La loi a posteriori est échantillonnée par Metropolis-Hastings à marche aléatoire, avec le diagnostic de Gelman et Rubin $\hat R$ quand plusieurs chaînes sont lancées.

\subsection{Juillet : modèle à grains 2D}

\subsubsection{Mise en données}

Une éprouvette 2D de $30 \times 60$~mm est découpée en grains de Voronoï de 3~mm, de volume élastique ($E = 52$~GPa, $\nu = 0{,}25$), et les joints de grain portent la moitié des résistances des joints intragranulaires. La résistance en traction des joints est fixée à 34~MPa, $G_I$ à 70~J/m$^2$ ; la cohésion et l'angle de frottement des joints sont libres sur $[4 ; 20]$~MPa et $[10 ; 35]^\circ$. Les cibles sont une traction directe de $18{,}3$~MPa et la résistance en compression simple, avec une dispersion de $10~\%$. Le plan compte 14 points et 3 réplicats ; deux émulateurs et quatre chaînes de 6\,000 itérations donnent la loi a posteriori.

\subsubsection{Résultats}

La loi a posteriori place l'angle de frottement à $13{,}4^\circ$ (intervalle à $95~\%$ de $10{,}2$ à $20{,}1^\circ$) et la cohésion à $13{,}6$~MPa ($6{,}9$ à $19{,}6$~MPa), avec $\hat R \leq 1{,}004$ (\cref{fig-posterior}). L'angle est collé à la borne basse de l'a priori : l'intervalle est en partie un effet de borne. Au jeu médian, sur trois tirages de grains, la traction vaut $17{,}07$~MPa ($-6{,}7~\%$) et la compression simple $139{,}1$~MPa ($+9{,}9~\%$).

\begin{figure}[htbp]
\centering
\includegraphics[width=0.75\linewidth]{calib/posterior_bohus}
\caption{Loi a posteriori de la calibration de juillet, modèle à grains 2D : angle de frottement des joints médian $13{,}4^\circ$, tronqué par la borne a priori ; cohésion peu contrainte.}
\label{fig-posterior}
\end{figure}

Hors calage, l'enveloppe triaxiale de ce jeu est trop plate (\cref{fig-triax-juillet}) : $\sigma_1 = 143{,}9 + 2{,}39\,\sigma_3$, soit un angle de frottement macroscopique de $24{,}2^\circ$ contre 45 à $55^\circ$ pour un granite, et $-56~\%$ à 20~MPa, $-62~\%$ à 50~MPa. Un balayage de l'angle des joints de $13{,}4$ à $45^\circ$ montre que la résistance simple croît trois fois plus vite que la pente de l'enveloppe : aucun angle ne cale les deux à la fois. Le confinement asservi reste en outre de 5 à $9~\%$ sous la consigne. Le jeu a été invalidé le 31 août, ses cibles mêmes étant contestées (traction de $18{,}3$ contre $10{,}3$~MPa, module de 52 contre $77{,}7$~GPa). Le 6 octobre, son rejeu avec le binaire courant a montré qu'il n'est plus reproductible (\cref{sec-calib-rejeu}).

\begin{figure}[htbp]
\centering
\includegraphics[width=0.8\linewidth]{calib/fig_triaxial}
\caption{Enveloppe du jeu de juillet. (a) Ajustement de Mohr-Coulomb des pics, angle macroscopique $24{,}2^\circ$. (b) Confinement atteint en fonction du confinement visé, inférieur de 5 à $9~\%$.}
\label{fig-triax-juillet}
\end{figure}

\subsubsection{Rejeu du 6 octobre 2026}
\label{sec-calib-rejeu}

Les quatre essais du jeu de juillet ont été rejoués le 6 octobre 2026 avec le binaire courant (\cle{f638b0f}, 4 fils). Les fichiers reprennent le jeu calé ; seuls la vitesse de chargement, sa rampe, la durée et le confinement changent d'un essai à l'autre (\cref{tab-calib-rejeu}). La traction tient : $16{,}67$~MPa contre $17{,}07$~MPa en juillet ($-2{,}3~\%$). La compression simple tombe à $33{,}03$~MPa contre $139{,}1$ et $141{,}5$~MPa archivés ($-76~\%$). Le jeu de juillet n'est donc plus reproductible.

\begin{table}[htbp]
\centering
\caption{Jeu de juillet rejoué le 6 octobre 2026 avec le binaire courant. Triaxiaux : pic brut, corrigé du biais des mors entre parenthèses ; écart à juillet sur le pic brut.}
\label{tab-calib-rejeu}
\small
\begin{tabular}{@{}lcccc@{}}
\toprule
essai & rejeu (MPa) & juillet (MPa) & écart à juillet & essai (MPa) \\
\midrule
traction directe & $16{,}67$ & $17{,}07$ & $-2{,}3~\%$ & $10{,}27$ (brésilien) \\
compression simple & $33{,}03$ & $139{,}1$ et $141{,}5$ & $-76~\%$ & $126{,}6$ \\
triaxial, $\sigma_3 = 20$~MPa & $69{,}22$ ($82{,}7$) & 186 & $-63~\%$ & $424{,}8$ \\
triaxial, $\sigma_3 = 50$~MPa & $122{,}01$ ($155{,}6$) & 247 & $-51~\%$ & $649{,}0$ \\
\bottomrule
\end{tabular}
\end{table}

Le même fichier de compression simple, à 1 fil, donne $33{,}04$~MPa avec le binaire du 25 août (\cle{139df12}) et avec celui du 19 août (\cle{5f8f71b}), premier commit de l'historique git. Le décrochage est donc antérieur à cet historique et ne peut pas être bissecté depuis le dépôt ; il faudrait les binaires \cle{rockim_f1} et \cle{rockim_f2} du poste local (\cref{tab-simus}). La fiche d'étude de rockim du 5 août 2026 consigne une chute de même ampleur, de $141{,}5$ à $32{,}8$~MPa, après la correction du signe de l'amortisseur de joint, qui anti-amortissait depuis l'écriture du code, et déclare déjà le jeu invalidé ; ce correctif est un candidat que la bissection départagerait.

La courbe de compression simple ne montre pas de rupture fragile. Le pic est atteint à $0{,}376$~ms sans aucun joint rompu ; un plateau à 32 à 33~MPa suit jusqu'à $0{,}54$~ms, puis une décroissance lente, avec 72 joints rompus sur $2\,054$ en fin de calcul. La résistance sature vers $0{,}97$ fois la résistance en traction des joints (34~MPa), avant toute rupture.

Le biais des mors bloqués est confirmé : à la fin de la rampe de confinement, la contrainte axiale vaut $6{,}56$ et $16{,}44$~MPa, à $-1{,}6$ et $-1{,}4~\%$ de $\nu/(1-\nu)\,\sigma_3$ ($\nu = 0{,}25$). Corrigés, les pics valent $82{,}7$ et $155{,}6$~MPa, très loin de juillet et de l'essai. Le confinement atteint au cœur est à $0{,}03~\%$ de la consigne, alors qu'il restait 5 à $9~\%$ dessous en juillet.

\subsection{Août : émulateur gaussien et prédiction triaxiale}

\subsubsection{Mise en données}

La deuxième campagne fige les choix numériques (insertion adaptative, adoucissement de Yan, contact par potentiel) et laisse libres six paramètres (\cref{tab-calib-aout}). L'éprouvette mesure $40 \times 80$~mm, avec le module mesuré de $77{,}7$~GPa ; la tessellation de Voronoï n'y sert que de maillage non structuré, toutes les phases étant identiques. Les cinq observables de la vraisemblance sont la résistance simple, le brésilien, $\sigma_1$ et la déformation au pic à 20~MPa, et la déformation au pic sans confinement. Les essais à 50, 75 et 100~MPa, hors calage, servent de prédiction. La base compte 60 calculs.

Le plafond déviatorique est un garde-fou du volume élastique : dans chaque élément, si la contrainte équivalente de von Mises dépasse le plafond, le déviateur est réduit à pression moyenne constante pour la ramener au plafond (\cref{sec-continu}).

\begin{table}[htbp]
\centering
\caption{Calibration d'août : lois a priori et a posteriori des six paramètres.}
\label{tab-calib-aout}
\small
\begin{tabular}{@{}lll@{}}
\toprule
paramètre & a priori & médiane et intervalle à $95~\%$ \\
\midrule
résistance en traction $f_t$ & $[20 ; 60]$~MPa & $25{,}7$ ($20{,}2$ à $43{,}9$)~MPa \\
cohésion $c$ & $[10 ; 90]$~MPa & $11{,}0$ ($10{,}0$ à $14{,}2$)~MPa, borne basse \\
angle de frottement & $[15 ; 60]^\circ$ & $46{,}2$ ($40{,}1$ à $49{,}8$)$^\circ$ \\
énergie $G_I$ & $[30 ; 300]$~J/m$^2$ & $50{,}1$ ($30{,}7$ à $115{,}7$)~J/m$^2$, borne basse \\
rapport $G_{II}/G_I$ & $[1 ; 12]$ & $2{,}63$ ($1{,}08$ à $8{,}41$) \\
plafond déviatorique & $[200 ; 1\,500]$~MPa & 462 (212 à 1\,242)~MPa \\
\bottomrule
\end{tabular}
\end{table}

\subsubsection{Résultats}

Avec le jeu retenu, calculé sur une fenêtre de temps assez longue pour atteindre le pic, $\sigma_1$ vaut $350{,}9$~MPa à 20~MPa ($-17~\%$), $814{,}4$~MPa à 50~MPa ($+25~\%$), puis $1\,002$ et $1\,222$~MPa à 75 et 100~MPa ($+29$ et $+36~\%$) : la prédiction échoue au-delà du critère de $20~\%$ (\cref{fig-calib-synthese}). L'enveloppe simulée est trop raide : elle croise l'enveloppe expérimentale vers 35~MPa, parce que l'enveloppe de Mohr-Coulomb linéaire des joints surestime les forts confinements. Un premier jeu voisin, calculé sur une fenêtre tronquée, semblait s'accorder à 50~MPa ($-3{,}6~\%$) : c'était un artefact de la troncature. La résistance simple et le brésilien n'ont pas été recalculés avec le jeu retenu ; le jeu voisin donnait $121{,}5$~MPa ($-4~\%$) et $12{,}25$~MPa ($+19~\%$).

Les déviateurs aux pics à 75 et 100~MPa, 927 et 1\,122~MPa, dépassent la médiane a posteriori du plafond déviatorique (462~MPa). Le jeu retenu porte un plafond de 897~MPa, que le déviateur à 100~MPa dépasse encore, ce qu'une moyenne de contraintes d'éléments plafonnées ne permet pas en norme de von Mises. L'écart, qui peut tenir à la mesure de la force aux plateaux ou à des effets dynamiques, n'a pas été instruit.

\begin{figure}[htbp]
\centering
\includegraphics[width=\linewidth]{calib/fig_calib_synthese}
\caption{Calibration Red Bohus. (a) Essais de Dumoulin et al. (b) Enveloppe : le jeu de juillet sous-estime le confinement, celui d'août le surestime au-delà d'environ 35~MPa. (c) Écart au pic expérimental.}
\label{fig-calib-synthese}
\end{figure}

L'émulateur a ensuite été infirmé par les calculs réels (coefficient de détermination croisé de $0{,}44$ sur la résistance simple) : son écart-type, 56~MPa, dépasse celui de l'essai, $21{,}4$~MPa, et c'est la taille du plan, 60 points en dimension 6, qui limite la calibration. L'échantillonnage n'a employé qu'une chaîne, sans $\hat R$, et deux paramètres sont collés à leur borne ; un échantillonneur à ensemble invariant affine triple la largeur de certains intervalles. Les essais confinés posaient le confinement avant la charge axiale avec les mors bloqués, ce qui peut biaiser le pic de $\nu/(1-\nu)\,\sigma_3$, soit 8 à 20~MPa ; ce biais, découvert en septembre, n'a pas été vérifié sur ces calculs ; il l'a été le 6 octobre sur le rejeu du jeu de juillet (\cref{sec-calib-rejeu}). Avec l'écart au plafond déviatorique, il laisse l'amplitude de l'échec à 75 et 100~MPa non consolidée.

\subsection{Septembre et calibrations continues}

Une troisième campagne, à 50~MPa de confinement, a testé une variable à la fois sur des calculs de moins de 5 minutes. Les écarts y portent sur le déviateur au pic. Corrigé du biais des mors, le déviateur au pic du cas homogène vaut 676~MPa, $13~\%$ au-dessus du déviateur mesuré de 599~MPa. Un jeu réduit à $0{,}7$ fois les résistances et à un module de 71~GPa donne $-5~\%$ à 50~MPa et $-9~\%$ à 20~MPa, et le même jeu au module de $77{,}7$~GPa $+10~\%$ à 100~MPa. Le seuil d'endommagement de dilatance reste collé au pic dans tous les cas ($0{,}95$ à $1{,}00$ du pic, contre $0{,}72$ mesuré). Les planchers numériques mesurés sont du même ordre que l'écart aux données : $1{,}1~\%$ pour le maillage, $3{,}7~\%$ pour la vitesse, 4 à $5~\%$ pour l'amortissement, $7~\%$ pour le schéma d'insertion. Une campagne de 150 calculs est préparée et n'a pas été lancée. La procédure publiée de Tatone et Grasselli~\cite{tatone2015}, qui cale successivement quatre familles de paramètres (volume, amortissement visqueux, rupture, interaction élastique) sur la compression simple, le brésilien et un essai biaxial, en est l'alternative.

Les calibrations par lois continues sous Abaqus éclairent le problème (\cref{fig-cdp}). Le modèle d'endommagement plastique calé à 50~MPa donne une enveloppe linéaire, $+25~\%$ à 20~MPa et $-7~\%$ à 100~MPa ; un Drucker-Prager hyperbolique donne une courbure de sens contraire à celle de l'essai. Trois représentations de l'hétérogénéité plafonnent à 310 à 357~MPa à 50~MPa de confinement, contre 650~MPa mesurés.

\begin{figure}[htbp]
\centering
\includegraphics[width=0.75\linewidth]{calib/cdp_envelope}
\caption{Abaqus, lois continues : le modèle d'endommagement plastique calé à 50~MPa donne une enveloppe linéaire, le Drucker-Prager hyperbolique une courbure de sens contraire.}
\label{fig-cdp}
\end{figure}

\subsection{Modèle à grains}

\subsubsection{Mise en données}

Le maillage de Voronoï découpe l'éprouvette en grains de phases minérales ; chaque propriété peut être redéfinie par phase, et les joints de grain prennent la moyenne des deux phases multipliée par un facteur propre à la frontière. Les résistances peuvent être tirées par joint selon une loi de Weibull de moyenne 1, avec un effet d'échelle identique à celui des VUMAT de la thèse :
\begin{equation}
P_f(\sigma) = 1 - \exp\left[-\left(\frac{\sigma}{\sigma_0}\right)^m\right], \qquad
f_t^J = f_t\left(\frac{Z_\text{eff}}{V_J}\right)^{1/m}, \quad V_J = \tfrac12\left(V_e^+ + V_e^-\right),\ Z_\text{eff} = 1~\text{mm}^3 .
\end{equation}
Avec l'exposant $m = 24$ retenu pour le Red Bohus, deux maillages dans le rapport de volume $3{,}16$ donnent des facteurs dans le rapport $3{,}16^{1/24}$ à $0{,}03~\%$ près.

\subsubsection{Vérifications}

Huit configurations d'éléments finis 3D vérifient le traitement des phases. Deux phases identiques donnent un résultat identique au bit près à celui d'une seule phase. Une barre de deux couches en série, de 83,1 et 29,3~GPa, a un module apparent de $44{,}16$~GPa pour $43{,}68$~GPa par la borne de Reuss,
\begin{equation}
\frac{1}{E_\text{Reuss}} = \sum_p \frac{f_p}{E_p},
\end{equation}
soit $1{,}1~\%$ d'écart. En compression élastique d'un cube de 1\,001 grains (fractions du tirage de grains : feldspath $62~\%$, quartz $31~\%$, biotite $7~\%$), le modèle à grains est $2{,}9~\%$ plus souple que le témoin homogène aux moyennes de Voigt. Le maximum de contrainte équivalente est $1{,}5$ fois plus fort ; la biotite porte en moyenne la moitié de la contrainte du témoin, le quartz $16~\%$ de plus.

En traction, un champ de Weibull corrélé à l'échelle du grain ($m = 6$) concentre la rupture dans la bande la plus faible (\cref{fig-weibull}) : le maillon faible rompt à $0{,}96\,f_t$, et trois maillages du même champ rompent dans la même bande à $4{,}6$ à $5{,}1$~MPa, contre $7{,}4$~MPa pour un champ non corrélé. Un banc bayésien synthétique retrouve un facteur de traction caché ($0{,}45$, intervalle $0{,}40$ à $0{,}46$) et laisse non identifié le facteur de cohésion, comme attendu en mode~I. Le modèle à grains a été abandonné le 2 septembre pour la calibration, faute de gain.

\begin{figure}[htbp]
\centering
\includegraphics[width=0.7\linewidth]{calib/fig_weibull}
\caption{Champ de Weibull corrélé sur les joints ($m = 6$, longueur de corrélation égale à la taille de grain) et fissure de traction finale, qui suit la bande la plus faible.}
\label{fig-weibull}
\end{figure}
